% 01-conventions.tex -- conventions and notation.
% Input by mostlyharmless.tex.  No preamble here; house macros are defined in
% the master file.

\section{Conventions and notation}
\label{sec:conventions}

This section fixes the conventions used everywhere else.  Later sections assume
it and do not restate it.  Where a symbol is overloaded in the wider literature,
the alternative reading is named explicitly, because those are the places where
errors enter silently.

\subsection{Frames and angles}

The observer frame is right-handed, with axes
\begin{equation}
  \hat e_x = \text{North}, \qquad
  \hat e_y = \text{East}, \qquad
  \hat e_z = \text{toward the observer},
  \label{eq:observer-frame}
\end{equation}
so that $\hat e_x \times \hat e_y = \hat e_z$, the third axis pointing from the sky
toward the observer --- into the ground for an observer looking up.  The line of
sight is $-\hat e_z$, and the disk is the unit disk seen along it.  A point on the
disk is specified either by Cartesian direction cosines $(x, y)$ or by plane polar
coordinates
\begin{equation}
  \rho = \sqrt{x^2 + y^2} = \sin\theta, \qquad
  \varphi = \arctan(y, x),
  \label{eq:disk-coords}
\end{equation}
where $\theta$ is the polar angle measured from $-\hat e_z$, and the
azimuth is measured from North through East, the same reference as the position
angle (D23).  For a source of unit angular radius, $\rho \in [0,1]$.

The body frame is the frame in which the surface map is specified, with $z$ along
the spin axis.  Its orientation relative to the observer is fixed by three
angles: the spin phase $\alpha$, the inclination $\beta$ measured from the line
of sight, and the position angle $\PA$ measured from North towards East.  The
body-to-observer rotation is
\begin{equation}
  R(\alpha,\beta,\PA) =
  R_z(\PA)\, R_y(\beta)\, R_z(\alpha),
  \label{eq:body-to-observer}
\end{equation}
with $R_y$ and $R_z$ active right-handed rotations.  The middle axis stays $y$:
the inclination is applied first, and the position angle then carries the result to
its place on the sky.  The body $z$ axis, which is the spin, maps to
$(\sin\beta\cos\PA,\ \sin\beta\sin\PA,\ \cos\beta)$, so it tilts toward the azimuth
$\PA$ and sits $\cos\beta$ above the sky plane.

\subsection{Spherical harmonics}

Complex harmonics are used in the derivations, and only for non-negative orders:
\begin{equation}
  Y_l^m(\theta,\varphi) = N_{lm}\, P_l^m(\cos\theta)\, e^{im\varphi},
  \qquad
  N_{lm} = \sqrt{\frac{2l+1}{4\pi}\,\frac{(l-m)!}{(l+m)!}} ,
  \label{eq:ylm}
\end{equation}
for $l \ge 0$ and $0 \le m \le l$, with the Condon--Shortley phase carried inside
$P_l^m$.  Negative orders are never formed --- not in the derivations, not for
maps, not anywhere --- so the magnitude $|m|$ is never needed either.  The order is
non-negative wherever it appears, and there is consequently no signed-order
convention to get wrong.

The harmonics are extended to negative orders by
\begin{equation}
  Y_l^{-m} := (-1)^m \, \overline{Y_l^{m}} , \qquad m > 0 ,
  \label{eq:ylm-extension}
\end{equation}
which is a definition, and which coincides with the standard convention built
from the associated Legendre function of negative order.  Negative orders are
wanted in exactly one place, the rotation machinery of section~3, where the order
is an index of the representation rather than of a map.

Coefficients, by contrast, are never formed at a negative order.  For a real map
the term accompanying $a_{lm} Y_l^m$ is $(-1)^m \overline{a_{lm}} \, Y_l^{-m}$,
which is to say $2\,\mathrm{Re}(a_{lm} Y_l^m)$; wherever the rotation would call
for a negative-order coefficient, its value is substituted in that form.

The basis used for maps is real.  At degree $l$ there are $2l+1$ real basis
functions, labelled by $m = 0, \dots, l$ together with a cosine or sine branch:
\begin{equation}
  R_l^0 = Y_l^0 , \qquad
  R_l^{m,c} = \sqrt{2}\,\mathrm{Re}\,Y_l^m , \qquad
  R_l^{m,s} = \sqrt{2}\,\mathrm{Im}\,Y_l^m \quad (m > 0) .
  \label{eq:real-basis}
\end{equation}
They are orthonormal on the sphere, and a real map is expanded in them with real
coefficients,
\begin{equation}
  I(\theta,\varphi) = \sum_{l} \Big[ \, c_{l0} R_l^0
  + \sum_{m=1}^{l} \big( c_{lm} R_l^{m,c} + s_{lm} R_l^{m,s} \big) \Big] .
  \label{eq:map-expansion}
\end{equation}

The constant term is $c_{00}R_0^0 = a_{00}Y_0^0 = a_{00}/\sqrt{4\pi}$, so $a_{00}$
is $\sqrt{4\pi}$ times the mean intensity over the sphere,
\begin{equation}
  \frac{1}{4\pi}\int_{S^2} I \, d\Omega = \frac{a_{00}}{\sqrt{4\pi}} .
  \label{eq:mean-intensity}
\end{equation}
It is neither the mean over the visible disk nor the flux.  Those integrals lack
the orthogonality of \eqref{eq:mean-intensity}: on the disk every harmonic picks up
the foreshortening factor $n$ of the visibility measure,
$\int_{\rm disk} Y_l^m \, d^2x = \int_{\rm visible} Y_l^m(\hat n)\, n \, d\Omega$,
which does not vanish for $l > 0$.

Maps in this work are normalised so that
\begin{equation}
  a_{00} = 1 ,
  \label{eq:a00-normalisation}
\end{equation}
so the uniform component is the reference and every other coefficient is a mode
amplitude relative to uniform brightness.  The sphere mean is then
$1/\sqrt{4\pi}$.  This is an assumption of the same kind as the unit radius: both
fix a scale that would otherwise be a free parameter, and both are stated once
rather than carried around.  Rotation preserves \eqref{eq:a00-normalisation},
because the $l = 0$ block is untouched by it --- the $d$ matrix at $l = 0$ is the
identity, which certificate \vtag{5} covers.  The free coefficients of a map are
therefore those with $l \ge 1$.

\begin{figure}[t]
  \centering
  \IfFileExists{generated/map.pdf}{%
    \includegraphics[width=0.8\textwidth]{generated/map.pdf}}{%
    \emph{Not generated yet --- run \code{make data}.}}
  \caption{A starspot on a uniformly bright surface, truncated at $l_{\max} = 16$,
  with brightness in units of the sphere mean that
  \eqref{eq:a00-normalisation} fixes.  Left: down the spin axis, near and far
  hemispheres.  Right: the same map rotated into the observer's frame by the
  rotation of section~\ref{sec:wigner}.  The ripples at the spot's rim are
  truncation error, the series being finite, and not structure.  The sky is drawn
  with North up and East to the left, and azimuths are measured from North through
  East (D23), so the spot's position on the page is quoted in that convention: a map
  referred to East instead would show the same spot a quarter turn away.}
  \label{fig:starspot}
\end{figure}

For compact storage the same map is carried as complex coefficients of the
non-negative orders alone,
\begin{equation}
  a_{l0} = c_{l0} , \qquad
  a_{lm} = \frac{c_{lm} - i\,s_{lm}}{\sqrt{2}} \quad (m > 0) ,
  \label{eq:half-complex}
\end{equation}
in terms of which \eqref{eq:map-expansion} reads
\begin{equation}
  I(\theta,\varphi) = \sum_{l} \Big[ \, a_{l0} Y_l^0
  + \sum_{m=1}^{l} 2\,\mathrm{Re}\big( a_{lm} Y_l^m \big) \Big] .
  \label{eq:map-from-half-complex}
\end{equation}
This is real by construction: there is no second set of stored coefficients that
could disagree with the first, and the explicit factor of $2$ is what a pair of
conjugate orders would otherwise have contributed.  At degree $\le L$ there are
$\sum_{l=0}^{L} (2l+1) = (L+1)^2$ independent real numbers, which is what a flat
coefficient vector holds.

\subsection{Visibility and Fourier conventions}

The complex visibility is the Fourier transform of the sky brightness, with the
factor $2\pi$ inside the exponent and no prefactor,
\begin{equation}
  V(u,v) = \int I(y,x)\, e^{-2\pi i (u x + v y)}\, dy\, dx ,
  \label{eq:visibility-definition}
\end{equation}
with $u$ and $v$ in \emph{wavelengths} and $y, x$ dimensionless direction
cosines.  Here $u$ is the \emph{North} component of the baseline and $v$ the
\emph{East} component: the reverse of the radio-interferometry convention, adopted
deliberately (D24) so that the coordinates on the sky and on the ground share an
orientation, and so that $\psi$ below is the position angle directly.  Thus
$V(0,0)$ is the total flux.  The transform is written in plane polar coordinates
through
\begin{equation}
  q = \sqrt{u^2 + v^2}, \qquad \psi = \arctan(v,u), \qquad a = 2\pi q .
  \label{eq:uv-polar}
\end{equation}
Public functions take $(u,v)$; $q$, $\psi$ and $a$ are internal.

The measure in \eqref{eq:visibility-definition} is the planar measure over the
projected disk, $dy\,dx = \rho\,d\rho\,d\varphi$.  The exact solid-angle
measure differs by a factor $1/n$ with $n = \sqrt{1-\rho^2}$, and is not used
here, because the planar measure is the one in which the classical uniform-disk
and limb-darkening results are expressed.  The difference is not cosmetic: the
solid-angle measure moves the square root from the odd family to the even one,
and so swaps which family is easy.

\subsection{Symbols}

\begin{center}
\begin{tabular}{ll}
\hline
symbol & meaning \\
\hline
$\theta,\ \varphi$ & polar and azimuthal angles on the sphere \\
$z = \cos\theta$ & natural variable for Legendre and Jacobi polynomials \\
$\rho = \sin\theta$ & disk radius; also $\sqrt{x^2+y^2}$ \\
$x,\ y$ & direction cosines: $x$ to North, $y$ to East \\
$u,\ v,\ w$ & baseline coordinates, in wavelengths \\
$q = \sqrt{u^2+v^2}$ & radial baseline, in wavelengths \\
$\psi = \arctan(v,u)$ & baseline azimuth in the $(u,v)$ plane \\
$a = 2\pi q$ & argument of the Bessel functions \\
$l,\ m$ & degree and order, with $m \ge 0$ throughout \\
$c_{lm},\ s_{lm}$ & real cosine and sine coefficients of a map \\
$a_{lm}$ & complex coefficient of order $m \ge 0$, \eqref{eq:half-complex} \\
$n$ & Zernike radial degree \\
$k = (n-m)/2$ & Jacobi degree, $m \ge 0$ \\
$\alpha,\ \beta,\ \PA$ & spin phase, inclination, position angle \\
$N_{lm}$ & harmonic normalisation, \eqref{eq:ylm} \\
$R_n^m(\rho)$ & Zernike radial polynomial, $m \ge 0$ \\
$j_l,\ J_m$ & spherical and cylindrical Bessel functions \\
\hline
\end{tabular}
\end{center}

Two overloadings are worth stating explicitly, because both have caused errors.
First, $u$ is a baseline coordinate here, whereas the limb-darkening literature
uses $u, u_1, u_2$ for the coefficients of the limb-darkening law; in this
package those coefficients are named \code{ld\_coeffs} and never $u$.  Second,
$y$ and $x$ are what the radio-interferometry literature calls the direction
cosines $l$ and $m$.  That notation is avoided here so that $l$ and $m$ always
mean harmonic degree and order.

\subsection{Verification conventions}

Claims that rest on algebra rather than on argument carry a tag such as
\vtag{11}, which resolves in the verification register to a named certificate
function in \code{mostlyharmless.symbolic}.  Each such claim also states its
scope and one of three evidence rungs:

\begin{description}
\item[exact] the residual simplifies to zero with the parameters left free;
\item[range-limited] the residual is exactly zero for every case in a stated
  finite set, typically all integer $l$ up to some $l_{\max}$ --- a proof of the
  finite statement, but not of the general one;
\item[numerical] agreement with an independent computation, to a stated
  tolerance.
\end{description}

The rung is not decoration.  ``Verified'' without a stated scope is the failure
mode this convention exists to prevent, and the distinction between a proof of a
finite statement and a proof of the general one is the one most easily blurred.
